\documentclass[11pt]{article}
\usepackage[T1]{fontenc}
\usepackage{lmodern}
\usepackage{amsmath,amssymb,amsthm}
\usepackage[margin=1in]{geometry}
\usepackage[colorlinks=true,linkcolor=blue,citecolor=blue,urlcolor=blue]{hyperref}
\usepackage{microtype}

\newtheorem{theorem}{Theorem}
\newtheorem{lemma}{Lemma}
\newcommand{\D}{\mathbb D}
\newcommand{\T}{\mathbb T}
\newcommand{\N}{\mathbb N}
\newcommand{\cD}{\mathcal D}

\title{Logarithmic cyclicity in Dirichlet-type spaces:\\
the complete phase diagram}
\author{Solution packet for arXiv:1301.4375, Problem 3}
\date{13 August 2026}

\begin{document}
\maketitle

\begin{abstract}
For
\[
 D_\alpha=\left\{f(z)=\sum_{n\geq0}a_nz^n:
 \sum_{n\geq0}(n+1)^\alpha|a_n|^2<\infty\right\},
\]
Problem 3 of B\'en\'eteau--Condori--Liaw--Seco--Sola asks whether
$f,\log f\in D_\alpha$ force $f$ to be cyclic.  We give the complete
answer: the implication holds precisely for $\alpha\geq0$.  The missing
positive range follows from the later complete-Pick theorem of
Aleman--Richter.  For every $\alpha<0$, a theorem of Dayan--Seco supplies a
singular measure on a Beurling--Carleson set whose Fourier coefficients,
after a weighted H\"older estimate, give a noncyclic singular inner function
$S_\mu$ with both $S_\mu$ and $\log S_\mu$ in $D_\alpha$.
\end{abstract}

\section{The question and the answer}

A function $f\in D_\alpha$ is \emph{cyclic} when the polynomial multiples of
$f$ are dense in $D_\alpha$, equivalently when there are polynomials $p_j$
such that $p_jf\to1$ in $D_\alpha$.  The phrase ``$\log f\in D_\alpha$''
means that $f$ has an analytic logarithm $q$ on $\D$ and that
$q\in D_\alpha$.

The source paper \cite{BCLSS} proves the implication for $\alpha>1$,
$\alpha=0$, and $\alpha=1$.  It obtains a conditional result for the
remaining parameters and asks whether its extra approximation hypothesis is
necessary.  The following settles the whole question.

\begin{theorem}[Complete phase diagram]\label{thm:main}
For a real number $\alpha$, the following are equivalent:
\begin{enumerate}
\item $\alpha\geq0$;
\item every $f$ satisfying $f,\log f\in D_\alpha$ is cyclic in $D_\alpha$.
\end{enumerate}
More precisely, if $\alpha<0$, there is a singular inner function $S_\mu$
such that $S_\mu,\log S_\mu\in D_\alpha$ but $S_\mu$ is not cyclic.
\end{theorem}

\paragraph{Proof intuition.}
There is a sharp change at the Hardy index.  On the nonnegative side,
$\log f\in D_\alpha\subset H^1$ forces $f$ to be outer; complete-Pick
Smirnov theory then turns the logarithm into a cyclicity certificate.  On the
negative side, $D_\alpha$ is large enough to contain every inner function.
We choose a singular inner function whose singular measure lives on a
Beurling--Carleson set, which forces noncyclicity, while arranging just enough
Fourier summability of that measure to put the analytic logarithm in
$D_\alpha$.

\section{The affirmative half}

We first recall the later theorem that closes the interval left open by the
source.

\begin{theorem}[Aleman--Richter \cite{AR}]\label{thm:AR}
Let $D(\nu)$ be a superharmonically weighted Dirichlet space.  If
$g\in D(\nu)$ is outer and
$\log g\in N^+(D(\nu))$, the Pick--Smirnov class of $D(\nu)$, then $g$ is
cyclic in $D(\nu)$.
\end{theorem}

The standard spaces $D_\alpha$, $0\leq\alpha\leq1$, belong to this family
\cite[Introduction]{AR}.  Their normalized complete Nevanlinna--Pick property
also gives
\begin{equation}\label{eq:smirnov}
 D_\alpha\subset N^+(D_\alpha).
\end{equation}

Suppose first that $0\leq\alpha\leq1$ and put $q=\log f$.  Since
$D_\alpha\subset H^2$, we have $q\in H^2\subset H^1$.  The standard
analytic-log criterion then says that $f$ is outer.  Indeed, radial boundary
values satisfy $\log|f^*|=\operatorname{Re}q^*$ almost everywhere, and the
$H^1$ mean-value formula gives
\[
 \log|f(0)|=\operatorname{Re}q(0)
 =\int_{\T}\operatorname{Re}q^*\,dm
 =\int_{\T}\log|f^*|\,dm,
\]
which is the outer equality (and $f=e^q$ has no zeros).  Now
$q\in D_\alpha\subset N^+(D_\alpha)$ by \eqref{eq:smirnov}, so
Theorem~\ref{thm:AR} makes $f$ cyclic.

For completeness, let $\alpha>1$.  The coefficient space $D_\alpha$ is a
Banach algebra: it is the one-sided weighted convolution space
$\ell^2_s$ with $s=\alpha/2>1/2$.  Hence the exponential power series
converges in $D_\alpha$, and
\[
 \frac1f=e^{-q}\in D_\alpha.
\]
Choose polynomials $p_j\to1/f$ in $D_\alpha$.  Multiplication by
$f\in D_\alpha$ is bounded, so $p_jf\to1$.  This proves cyclicity for all
$\alpha\geq0$.

\section{Counterexamples for every negative parameter}

We use two facts assembled explicitly in Dayan--Seco \cite{DS}.  First, the
Korenblum--Roberts theorem says that, for $\alpha<0$, a singular inner
function $S_\mu$ is cyclic in $D_\alpha$ if and only if
\begin{equation}\label{eq:KR}
 \mu(E)=0\quad\hbox{for every Beurling--Carleson set }E\subset\T.
\end{equation}
Second, their Theorem 6.3 states that for each $p>2$ there is a nonzero
positive singular measure $\mu$, supported on a Beurling--Carleson set, such
that
\begin{equation}\label{eq:lp}
 (\widehat\mu(n))_{n\in\mathbb Z}\in\ell^p(\mathbb Z).
\end{equation}
Their proof uses a Salem measure and also verifies the
Beurling--Carleson property of its support.

Fix $\alpha<0$.  Select $p>2$ with
\begin{equation}\label{eq:pchoice}
 \frac{\alpha p}{p-2}<-1.
\end{equation}
If $\alpha\leq-1$, every $p>2$ works.  If $-1<\alpha<0$, it is enough to
take
\[
 2<p<\frac{2}{1+\alpha},
\]
an interval which is nonempty precisely because $\alpha<0$.
Take the measure supplied by \eqref{eq:lp} for this $p$.  H\"older's
inequality with conjugate exponents $p/2$ and $p/(p-2)$ yields
\begin{align}
 \sum_{n\geq1}(n+1)^\alpha|\widehat\mu(n)|^2
 &\leq
 \left(\sum_{n\geq1}|\widehat\mu(n)|^p\right)^{2/p}
 \left(\sum_{n\geq1}(n+1)^{\alpha p/(p-2)}\right)^{(p-2)/p}
 <\infty,                                      \label{eq:holder}
\end{align}
where the final series converges by \eqref{eq:pchoice}.

Define the associated singular inner function
\[
 S_\mu(z)=\exp\left\{-\int_{\T}\frac{\zeta+z}{\zeta-z}\,d\mu(\zeta)\right\}.
\]
Its defining exponent is an analytic logarithm.  Expanding the Herglotz
kernel gives, up to the choice of Fourier-sign convention,
\begin{equation}\label{eq:logexpansion}
 \log S_\mu(z)=-\mu(\T)-2\sum_{n\geq1}\widehat\mu(n)z^n.
\end{equation}
Thus \eqref{eq:holder} says that $\log S_\mu\in D_\alpha$.  Moreover,
$S_\mu\in H^\infty\subset H^2$, and $(n+1)^\alpha\leq1$ shows that
$H^2\subset D_\alpha$.  Consequently $S_\mu\in D_\alpha$ as well.

Finally, $\mu$ is supported on a Beurling--Carleson set $E$ and is nonzero,
so $\mu(E)=\mu(\T)>0$.  Condition \eqref{eq:KR} fails, and $S_\mu$ is not
cyclic in $D_\alpha$.  This proves the negative half and completes the proof
of Theorem~\ref{thm:main}. \qed

\section{Audit of the upgrade and its scope}

The atomic singular inner function
$\exp(-(1+z)/(1-z))$ already gives counterexamples for $\alpha<-1$, because
the coefficients of its logarithm have constant magnitude.  It does not
reach $-1\leq\alpha<0$.  The substantive upgrade is the choice of $p$ in
\eqref{eq:pchoice}, which converts Dayan--Seco's $\ell^p$ Fourier measure into
the exact weighted $\ell^2$ regularity required for every negative
$\alpha$.

As a cross-check, the Salem estimate quoted in their proof is
$\widehat\mu(n)=O(n^{-s/2+\varepsilon})$ for arbitrary $s\in(0,1)$, on a
support they show to be Beurling--Carleson.  Choosing
$s-2\varepsilon>1+\alpha$ makes the series in \eqref{eq:holder} converge
directly.  This gives the same endpoint coverage without the abstract
embedding calculation.

This packet is a literature-assisted synthesis.  Aleman--Richter state the
theorem used for the positive range.  Dayan--Seco state the measure theorem
and the Korenblum--Roberts criterion, but do not state the complete
$D_\alpha$ phase diagram verbatim.  Their 2026 publisher correction only
repairs a typeset integral in another part of the proof and does not affect
Theorem 6.3 or the argument above \cite{DScorr}.

\begin{thebibliography}{9}

\bibitem{BCLSS}
C. B\'en\'eteau, A. A. Condori, C. Liaw, D. Seco, and A. A. Sola,
\emph{Cyclicity in Dirichlet-type spaces and extremal polynomials},
J. Anal. Math. \textbf{126} (2015), 259--286;
\href{https://arxiv.org/abs/1301.4375}{arXiv:1301.4375}.

\bibitem{AR}
A. Aleman and S. Richter,
\emph{Cyclicity and iterated logarithms in the Dirichlet space},
Proc. Amer. Math. Soc. \textbf{154} (2026), no.~2, 783--791;
\href{https://arxiv.org/abs/2409.20298}{arXiv:2409.20298}.

\bibitem{DS}
A. Dayan and D. Seco,
\emph{On the cyclic behavior of singular inner functions in Besov and
sequence spaces},
Banach J. Math. Anal. \textbf{20} (2026), article 2;
\href{https://doi.org/10.1007/s43037-025-00467-w}{doi:10.1007/s43037-025-00467-w},
\href{https://arxiv.org/abs/2504.05208}{arXiv:2504.05208}.

\bibitem{DScorr}
A. Dayan and D. Seco,
\emph{Correction: On the cyclic behavior of singular inner functions in
Besov and sequence spaces},
Banach J. Math. Anal. \textbf{20} (2026), article 25;
\href{https://doi.org/10.1007/s43037-026-00488-z}{doi:10.1007/s43037-026-00488-z}.

\bibitem{K}
B. Korenblum,
\emph{Cyclic elements in some spaces of analytic functions},
Bull. Amer. Math. Soc. \textbf{5} (1981), no.~3, 317--318.

\bibitem{R}
J. W. Roberts,
\emph{Cyclic inner functions in the Bergman spaces and weak outer functions
in $H^p$, $0<p<1$},
Illinois J. Math. \textbf{29} (1985), 25--38.

\bibitem{S}
R. Salem,
\emph{On singular monotonic functions whose spectrum has a given Hausdorff
dimension}, Ark. Mat. \textbf{1} (1951), 353--365.

\end{thebibliography}

\end{document}
